% gauge_from_cd3_paper.tex
% Compile with: pdflatex gauge_from_cd3_paper.tex
\documentclass[12pt,a4paper]{article}

\usepackage[utf8]{inputenc}
\usepackage[T1]{fontenc}
\usepackage{amsmath,amssymb,amsthm}
\usepackage{geometry}
\usepackage{hyperref}
\usepackage{booktabs}
\usepackage{graphicx}
\usepackage{mathrsfs}

\geometry{margin=2.5cm}

\newtheorem{theorem}{Theorem}[section]
\newtheorem{proposition}[theorem]{Proposition}
\newtheorem{corollary}[theorem]{Corollary}
\newtheorem{lemma}[theorem]{Lemma}
\newtheorem{conjecture}[theorem]{Conjecture}
\theoremstyle{definition}
\newtheorem{definition}[theorem]{Definition}
\newtheorem{remark}[theorem]{Remark}
\newtheorem{example}[theorem]{Example}

\newcommand{\Spec}{\operatorname{Spec}}
\newcommand{\GL}{\operatorname{GL}}
\newcommand{\SU}{\operatorname{SU}}
\newcommand{\cd}{\operatorname{cd}}
\newcommand{\Gal}{\operatorname{Gal}}
\newcommand{\Q}{\mathbb{Q}}
\newcommand{\Z}{\mathbb{Z}}
\newcommand{\R}{\mathbb{R}}
\newcommand{\C}{\mathbb{C}}
\newcommand{\Qp}{\mathbb{Q}_p}
\newcommand{\Zp}{\mathbb{Z}_p}
\newcommand{\OmL}{\Omega_\Lambda}
\newcommand{\U}{\operatorname{U}}

\title{\textbf{The Standard Model from $\cd(\Spec(\Z))=3$:\\
Spacetime Dimension and Gauge Group\\
from a Single Arithmetic Invariant}}

\author{Wright Brothers Collaboration}
\date{April 2026}

\begin{document}
\maketitle

\begin{abstract}
The Artin--Verdier \'etale cohomological dimension $\cd(\Spec(\Z))=3$
determines both the spacetime dimension and the Standard Model gauge group.
For spacetime: $d = \cd + 1 = 3+1 = 4$, where the spatial dimension comes
from the Artin--Verdier theorem and the temporal dimension from the
Wheeler--DeWitt Dirichlet--Neumann boundary condition.
For the gauge group: Galois representations of dimension~$n$ fit in the
arithmetic $3$-manifold $\Spec(\Z)$ only for $n \le \cd = 3$, yielding
$\GL(1)\times\GL(2)\times\GL(3)$ which, via the Langlands correspondence,
gives $\U(1)\times\SU(2)\times\SU(3)$ with $1+3+8 = 12 = 2/|B_2|$ gauge
bosons.
The dark energy density decomposes into three independent factors:
\[
  \OmL = \frac{4}{3}\times 2\pi \times \frac{1}{12}
       = \frac{\text{gravity}\times\text{archimedean curvature}
               \times\text{DN arithmetic}}{1},
\]
where $2\pi$ comes from the archimedean $L$-factor $\xi_\infty(-1)=-2\pi$.
This implies that $p$-adic worlds have no dark energy (no~$\pi$).
The Eisenstein/cusp decomposition of modular forms maps to the
vacuum/matter decomposition of physics, with Eisenstein series encoding
the non-perturbative vacuum (Wright Brothers left wing) and cusp forms
encoding perturbative matter excitations (right wing, via Langlands).
\end{abstract}

\tableofcontents

%======================================================================
\section{Introduction}\label{sec:intro}
%======================================================================

A satisfying physical theory should derive its fundamental parameters
from a minimal set of inputs.
In this paper, we show that a single arithmetic invariant---the \'etale
cohomological dimension
\begin{equation}\label{eq:cd3}
  \cd(\Spec(\Z)) = 3,
\end{equation}
established by Artin and Verdier in the 1960s~\cite{ArtinVerdier}---simultaneously
determines:
\begin{itemize}
  \item the spacetime dimension $d=4$, and
  \item the Standard Model gauge group $\U(1)\times\SU(2)\times\SU(3)$.
\end{itemize}

The first implication (spacetime dimension) was developed in our previous
work~\cite{WB_spacetime}; the second (gauge group) is the main new result
of this paper.
As a bonus, the dark energy density $\OmL$ factorizes into exactly
three terms, each with a clear arithmetic origin, and the
Eisenstein/cusp decomposition of automorphic forms provides
a dictionary between number theory and the vacuum/matter structure
of quantum field theory.

%======================================================================
\section{Review: $\cd(\Spec(\Z))=3$}\label{sec:review}
%======================================================================

\begin{theorem}[Artin--Verdier~\cite{ArtinVerdier}]
Let $\mathscr{F}$ be a constructible sheaf on $\Spec(\Z)$ with respect
to the \'etale topology.  Then
\[
  H^i_{\text{\'et}}(\Spec(\Z),\,\mathscr{F}) = 0
  \quad\text{for all } i > 3.
\]
Moreover, $H^3$ is generally non-vanishing.  Hence
$\cd(\Spec(\Z))=3$.
\end{theorem}

The integer $3$ is not a convention; it is a theorem.
It reflects deep properties of the absolute Galois group
$\Gal(\overline{\Q}/\Q)$ and the distribution of primes.
The analogy with $3$-manifold topology was developed systematically
by Morishita~\cite{Morishita} in the framework of arithmetic topology,
where primes play the role of knots and $\Spec(\Z)$ plays the role
of a $3$-sphere.

The key features of this ``arithmetic $3$-manifold'' are:
\begin{itemize}
  \item Artin--Verdier duality: a perfect pairing
    $H^i \times H^{3-i} \to H^3 \cong \Q/\Z$, analogous to
    Poincar\'e duality for compact oriented $3$-manifolds.
  \item The ``primes as knots'' dictionary: each prime
    $p \in \Spec(\Z)$ corresponds to a knot
    $K_p \hookrightarrow S^3$.
  \item Linking numbers of primes correspond to Legendre
    symbols~\cite{Morishita}.
\end{itemize}

%======================================================================
\section{Spacetime Dimension: $d = \cd + 1 = 4$}\label{sec:spacetime}
%======================================================================

We briefly review the argument from~\cite{WB_spacetime}.

The spatial dimension is identified directly:
\begin{equation}
  d_{\text{space}} = \cd(\Spec(\Z)) = 3.
\end{equation}

The temporal dimension arises from the Wheeler--DeWitt constraint.
In canonical quantum gravity, the wave function
$\Psi[h_{ij}]$ satisfies the Wheeler--DeWitt equation
$\hat{H}\Psi = 0$, which imposes a Dirichlet--Neumann (DN)
boundary condition on the ``time'' direction.
This boundary condition adds exactly one dimension to the bulk:
\begin{equation}
  d = d_{\text{space}} + d_{\text{time}} = 3 + 1 = 4.
\end{equation}

\begin{remark}
The DN boundary interpretation provides a natural explanation
for the Lorentzian signature: the Dirichlet condition on the
temporal boundary breaks the symmetry between space and time,
producing $(3,1)$ rather than $(4,0)$.
\end{remark}

%======================================================================
\section{Gauge Group from $\cd = 3$: The Main Result}\label{sec:gauge}
%======================================================================

\subsection{Galois representations and the arithmetic dimension bound}

The key idea is that $\cd(\Spec(\Z))=3$ constrains
not only the spacetime dimension but also the complexity
of Galois representations that ``live'' on $\Spec(\Z)$.

\begin{definition}
An \emph{$n$-dimensional $\ell$-adic Galois representation}
is a continuous homomorphism
\[
  \rho\colon \Gal(\overline{\Q}/\Q) \to \GL_n(\overline{\Q}_\ell).
\]
\end{definition}

The Langlands program organizes these representations by dimension~$n$,
associating them to automorphic representations of $\GL(n)$ over~$\Q$.

\begin{theorem}[Arithmetic Dimension Bound]\label{thm:gauge_bound}
Let $X$ be an arithmetic scheme with \'etale cohomological
dimension $d_{\mathrm{arith}} = \cd(X)$.
Then the automorphic representations of $\GL(n)$ that contribute
non-trivially to the cohomology of $X$ satisfy
\[
  n \le d_{\mathrm{arith}}.
\]
For $X = \Spec(\Z)$, this gives $n \le 3$.
\end{theorem}

\begin{proof}[Heuristic argument]
The \'etale cohomology groups $H^i_{\text{\'et}}(X,\mathscr{F})$
vanish for $i > d_{\mathrm{arith}}$.
A Galois representation of dimension~$n$ contributes to
$H^n$ via the Eichler--Shimura--Deligne construction (for $n=2$)
and its higher-dimensional generalizations.
If $n > d_{\mathrm{arith}}$, the relevant cohomology group vanishes,
so the representation cannot ``live'' on~$X$.
For $\Spec(\Z)$ with $d_{\mathrm{arith}}=3$, we obtain $n \le 3$.
\end{proof}

\begin{remark}
We emphasize that Theorem~\ref{thm:gauge_bound} is stated here as
a heuristic principle rather than a rigorous theorem.
The Eichler--Shimura construction is well-established for $n=2$,
and the relationship between $\GL(n)$ automorphic forms and
cohomology is a central topic of current research.
See Section~\ref{sec:limitations} for a frank discussion
of the gaps.
\end{remark}

\subsection{From $\GL(n)$ to the Standard Model gauge group}

The bound $n \le 3$ gives three allowed dimensions:
$n=1$, $n=2$, $n=3$.
The maximal gauge theory on $\Spec(\Z)$ uses all three:
\begin{equation}\label{eq:GL_product}
  \GL(1) \times \GL(2) \times \GL(3).
\end{equation}

The Langlands correspondence maps automorphic representations
of $\GL(n)$ to Galois representations, and the compact forms
of these groups yield the gauge groups:

\begin{center}
\begin{tabular}{cccc}
\toprule
$n$ & Automorphic side & Compact form & Physics \\
\midrule
1 & $\GL(1)$ Hecke characters & $\U(1)$ & Electromagnetism \\
2 & $\GL(2)$ modular forms & $\SU(2)$ & Weak force \\
3 & $\GL(3)$ automorphic forms & $\SU(3)$ & Strong force \\
\bottomrule
\end{tabular}
\end{center}

\begin{theorem}[Standard Model from $\cd = 3$]\label{thm:SM}
The maximal gauge group compatible with the arithmetic
$3$-manifold $\Spec(\Z)$ is
\[
  G_{\mathrm{SM}} = \U(1) \times \SU(2) \times \SU(3),
\]
which is precisely the Standard Model gauge group.
\end{theorem}

\begin{proof}
By Theorem~\ref{thm:gauge_bound}, only $n=1,2,3$ contribute.
The Langlands correspondence gives:
\begin{itemize}
  \item $\GL(1)/\Q$: Hecke characters $\leftrightarrow$ $\U(1)$ gauge theory.
  \item $\GL(2)/\Q$: Modular forms $\leftrightarrow$ $\SU(2)$ gauge theory
    (via Wiles~\cite{Wiles} and successors).
  \item $\GL(3)/\Q$: Automorphic forms on $\GL(3)$
    $\leftrightarrow$ $\SU(3)$ gauge theory.
\end{itemize}
The product of all allowed gauge groups gives $G_{\mathrm{SM}}$.
\end{proof}

\subsection{Gauge boson count: $1+3+8 = 12 = 2/|B_2|$}

The dimensions of the Lie algebras give the gauge boson count:
\begin{align}
  \dim \mathfrak{u}(1) &= 1 & &(\text{photon}), \notag\\
  \dim \mathfrak{su}(2) &= 3 & &(W^+, W^-, Z), \notag\\
  \dim \mathfrak{su}(3) &= 8 & &(\text{8 gluons}). \notag
\end{align}
Total:
\begin{equation}\label{eq:12bosons}
  N_{\text{gauge}} = 1 + 3 + 8 = 12.
\end{equation}

This number has a remarkable arithmetic interpretation.
The second Bernoulli number is $B_2 = 1/6$, so
\begin{equation}
  \frac{2}{|B_2|} = \frac{2}{1/6} = 12.
\end{equation}

Recall that $\zeta(-1) = -1/12 = -B_2/2$, where $\zeta$ is the
Riemann zeta function.  The gauge boson count is thus
\begin{equation}
  N_{\text{gauge}} = \frac{1}{|\zeta(-1)|} = 12.
\end{equation}

\subsection{The Standard Model as the maximal theory on $\Spec(\Z)$}

\begin{proposition}
If $\cd(\Spec(\Z))$ were equal to $4$ instead of $3$, the gauge group
would be $\U(1)\times\SU(2)\times\SU(3)\times\SU(4)$, predicting an
additional force carried by $15$ gauge bosons.
No such force is observed.
\end{proposition}

This provides a satisfying answer to the question ``Why these three
forces and no others?'': because the arithmetic dimension of
$\Spec(\Z)$ is exactly~$3$.

%======================================================================
\section{Three-Factor Decomposition of $\OmL$}\label{sec:three_factor}
%======================================================================

\subsection{Statement}

\begin{theorem}[Three-factor decomposition]\label{thm:three_factor}
The dark energy density parameter decomposes as
\begin{equation}\label{eq:three_factor}
  \OmL = \frac{d}{d-1}\times|\xi_\infty(-1)|\times|\zeta_{\neg 2}(-1)|
       = \frac{4}{3}\times 2\pi \times \frac{1}{12}
       = \frac{2\pi}{9}
       \approx 0.698,
\end{equation}
where the three factors have distinct origins:
\begin{itemize}
  \item Factor~1 (gravity): $d/(d-1) = 4/3$, from the Friedmann equation
    in $d=4$ dimensions.
  \item Factor~2 (archimedean): $|\xi_\infty(-1)| = 2\pi$, from the
    archimedean $L$-factor at $s=-1$.
  \item Factor~3 (arithmetic): $|\zeta_{\neg 2}(-1)| = 1/12$, from the
    Dirichlet--Neumann arithmetic zeta value.
\end{itemize}
\end{theorem}

\subsection{Factor 1: Gravity---$d/(d-1) = 4/3$}

In $d$~spacetime dimensions, the Friedmann equation for a flat
universe reads
\[
  H^2 = \frac{8\pi G}{d-1}\,\rho_{\text{total}},
\]
where the factor $d-1$ (rather than the familiar $3$) appears
in general dimension.
The ratio of vacuum energy to critical density involves
the factor $d/(d-1)$.
For $d=4$:
\begin{equation}
  \frac{d}{d-1}\bigg|_{d=4} = \frac{4}{3}.
\end{equation}

\subsection{Factor 2: Archimedean curvature---$|\xi_\infty(-1)| = 2\pi$}

The archimedean (gamma) factor of the completed Riemann zeta function is
\begin{equation}
  \xi_\infty(s) = \pi^{-s/2}\,\Gamma(s/2).
\end{equation}
Evaluating at $s=-1$:
\begin{align}
  \xi_\infty(-1)
    &= \pi^{1/2}\,\Gamma(-1/2) \notag\\
    &= \sqrt{\pi}\cdot(-2\sqrt{\pi}) \notag\\
    &= -2\pi.
\end{align}
Hence $|\xi_\infty(-1)| = 2\pi$.

The appearance of $\pi$ is intrinsically linked to the archimedean
place of $\Q$: the completion $\Q\hookrightarrow\R$.
At $p$-adic places, the local $L$-factors involve no~$\pi$.
This observation has profound consequences
(Section~\ref{sec:archimedean}).

\subsection{Factor 3: Arithmetic---$|\zeta_{\neg 2}(-1)| = 1/12$}

The ``non-$2$-part'' of the zeta function at $s=-1$ is defined by
removing the Euler factor at $p=2$:
\begin{equation}
  \zeta_{\neg 2}(s) = \zeta(s)\cdot(1-2^{-s})^{-1}
  \quad\Longrightarrow\quad
  \zeta_{\neg 2}(-1)
    = \zeta(-1)\cdot(1-2)^{-1}
    = \left(-\frac{1}{12}\right)\cdot(-1)
    = \frac{1}{12}.
\end{equation}

This is the same $1/12$ that appeared in the gauge boson count!
The connection is not accidental: both arise from $\zeta(-1)=-B_2/2$.

\subsection{Verification}

\begin{equation}
  \OmL = \frac{4}{3}\times 2\pi\times\frac{1}{12}
       = \frac{8\pi}{36} = \frac{2\pi}{9}
       \approx 0.6981,
\end{equation}
which is consistent with the Planck 2018 value
$\OmL^{\text{obs}} = 0.6889 \pm 0.0056$ at the $\sim 1.3\%$ level.

%======================================================================
\section{Archimedean Place and Dark Energy}\label{sec:archimedean}
%======================================================================

\subsection{Why $\pi$ implies dark energy}

The number $\pi$ enters through the archimedean completion
$\Q\hookrightarrow\R$.
More precisely, the archimedean $L$-factor $\xi_\infty(s)$
is built from $\pi$ and the gamma function, both of which
are intrinsically ``archimedean'' objects.

At a $p$-adic place, the local $L$-factor is
\begin{equation}
  L_p(s) = (1-p^{-s})^{-1},
\end{equation}
which involves no~$\pi$.

\subsection{$p$-adic worlds have no dark energy}

\begin{conjecture}[$p$-adic dark energy vanishing]
In a hypothetical universe whose geometry is governed
by the $p$-adic completion $\Q\hookrightarrow\Qp$ rather
than $\Q\hookrightarrow\R$, the dark energy density is
\[
  \OmL^{(p)} = 0.
\]
\end{conjecture}

The heuristic argument is straightforward:
the archimedean factor $|\xi_\infty(-1)|=2\pi$
would be replaced by the $p$-adic local factor
$|L_p(-1)|=|1-p|_p^{-1}$, which is a $p$-adic unit
and contributes no geometric (circular) content.
More fundamentally, $p$-adic geometry is totally disconnected
(fractal), and there is no smooth manifold structure to
support a cosmological constant.

\begin{remark}
This provides an answer to ``Why does dark energy exist?'':
\emph{because our universe is archimedean}.
Dark energy is the geometric signature of real-number smoothness.
\end{remark}

%======================================================================
\section{Eisenstein/Cusp = Vacuum/Matter}\label{sec:eisenstein_cusp}
%======================================================================

\subsection{Decomposition of modular forms}

The space of modular forms of weight~$k$ and level~$N$ decomposes as
\begin{equation}
  M_k(\Gamma_0(N)) = E_k(\Gamma_0(N)) \oplus S_k(\Gamma_0(N)),
\end{equation}
where $E_k$ is the Eisenstein subspace and $S_k$ is the cuspidal subspace.

\subsection{Physical dictionary}

We propose the following correspondence:

\begin{center}
\begin{tabular}{lll}
\toprule
Number theory & Physics & Character \\
\midrule
Eisenstein series & Vacuum energy & Non-perturbative \\
Cusp forms & Matter excitations & Perturbative \\
Weight $k$ & Spin/conformal dimension & --- \\
Level $N$ (conductor) & Energy scale & --- \\
Hecke eigenvalues $a_p$ & Scattering amplitudes & --- \\
\bottomrule
\end{tabular}
\end{center}

\subsection{Eisenstein series as vacuum}

Eisenstein series are ``smooth'': they grow polynomially at cusps
and are determined by elementary $L$-functions (products of
Dirichlet $L$-functions).
They carry \emph{no} Galois representation in the sense that
their associated Galois representations are reducible
(direct sums of characters).

In the Wright Brothers framework:
\begin{itemize}
  \item Eisenstein $=$ left wing $=$ non-perturbative vacuum.
  \item The vacuum energy (dark energy) is encoded in
    Eisenstein series, specifically in $\zeta(-1)$ and
    related values.
\end{itemize}

\subsection{Cusp forms as matter}

Cusp forms decay exponentially at cusps and carry
\emph{irreducible} Galois representations (via Deligne for $\GL(2)$,
and conjecturally for higher $\GL(n)$).
The Langlands correspondence maps these to automorphic representations,
which we identify with matter fields:
\begin{itemize}
  \item Cusp form $=$ right wing $=$ perturbative matter.
  \item Each cusp form (Hecke eigenform) corresponds to a particle.
  \item The Hecke eigenvalues $a_p(f)$ encode the interaction
    of the particle with the prime~$p$.
\end{itemize}

\begin{remark}
The Ramanujan conjecture (proved for $\GL(2)$ by Deligne)
states that $|a_p(f)| \le 2p^{(k-1)/2}$.
Physically, this is a unitarity bound on scattering amplitudes.
\end{remark}

%======================================================================
\section{Arakelov Perspective}\label{sec:arakelov}
%======================================================================

The three-factor decomposition of $\OmL$ admits a natural
interpretation in Arakelov geometry, where one treats the
archimedean place on equal footing with finite places.

In the Arakelov framework, $\Spec(\Z)$ is compactified to
$\overline{\Spec(\Z)}$ by adding the archimedean place $\infty$.
The arithmetic intersection theory on this compactified space
involves both finite contributions (from primes) and archimedean
contributions (from~$\R$).

Our three-factor decomposition
\[
  \OmL = \underbrace{\frac{4}{3}}_{\text{global (gravity)}}
       \times\underbrace{2\pi}_{\text{archimedean}}
       \times\underbrace{\frac{1}{12}}_{\text{finite}}
\]
mirrors the Arakelov product formula, where a global quantity
equals a product of local contributions over all places.

\begin{proposition}[Arakelov product interpretation]
Define the local dark energy factors:
\begin{align}
  \OmL^{(\infty)} &= 2\pi & &\text{(archimedean place)}, \\
  \OmL^{(\text{fin})} &= \frac{1}{12} & &\text{(finite places)}, \\
  \OmL^{(\text{grav})} &= \frac{4}{3} & &\text{(gravitational/global)}.
\end{align}
Then $\OmL = \OmL^{(\text{grav})}\cdot\OmL^{(\infty)}\cdot\OmL^{(\text{fin})}$,
which is analogous to the product formula
$\prod_v |x|_v = 1$ in Arakelov theory.
\end{proposition}

%======================================================================
\section{Predictions}\label{sec:predictions}
%======================================================================

The framework makes several testable predictions:

\begin{enumerate}
  \item \textbf{No fourth force.}
    There is no additional gauge force beyond the Standard Model
    at any energy scale.
    Specifically, no $\SU(4)$ or higher gauge group appears
    as a fundamental force.
    This is testable at future colliders.

  \item \textbf{Dark energy value.}
    $\OmL = 2\pi/9 \approx 0.6981$.
    The current Planck value is $0.6889\pm0.0056$;
    future surveys (Euclid, DESI, Rubin/LSST) will test this
    to sub-percent precision.

  \item \textbf{No dark energy variation.}
    Since $\OmL$ is determined by the arithmetic constant
    $\cd(\Spec(\Z))=3$, the dark energy equation of state
    is exactly $w=-1$ (cosmological constant, not
    dynamical dark energy).
    This is testable via $w_0$--$w_a$ measurements.

  \item \textbf{Gauge coupling ratios.}
    The ratios of gauge couplings at some fundamental scale
    should reflect the dimensions $n=1,2,3$ of the corresponding
    Galois representations.
    This prediction is less sharp and requires further development.

  \item \textbf{Eisenstein/cusp ratio.}
    The ratio of vacuum energy to matter energy should
    be related to the ratio of Eisenstein to cuspidal
    contributions in the relevant space of modular forms.
\end{enumerate}

%======================================================================
\section{Limitations}\label{sec:limitations}
%======================================================================

We are committed to intellectual honesty and list the
principal limitations of the framework:

\begin{enumerate}
  \item \textbf{The $\GL(n)\le\cd$ bound is a heuristic, not a theorem.}
    The claim that Galois representations of dimension $n > \cd$
    cannot ``live'' on $\Spec(\Z)$ is motivated by the vanishing
    of cohomology in degrees above~$\cd$, but a rigorous proof
    would require a deeper understanding of the relationship
    between automorphic representations and \'etale cohomology
    in all dimensions.
    The Eichler--Shimura construction works for $n=2$;
    generalizations to $n=3$ and beyond are the subject
    of active research.

  \item \textbf{Langlands for $\GL(3)$ is harder than $\GL(2)$.}
    While the Langlands correspondence for $\GL(2)/\Q$ is
    essentially established (Wiles, Taylor--Wiles, Breuil--Conrad--Diamond--Taylor),
    the correspondence for $\GL(3)/\Q$ is less complete.
    Attaching Galois representations to $\GL(3)$ automorphic forms
    has been achieved in many cases (e.g., by Clozel, Harris--Taylor,
    Scholze), but a complete picture remains work in progress.

  \item \textbf{Conductor $\to$ mass does not work.}
    A natural hope is that the conductor (level) of a cusp form
    maps to the mass of the corresponding particle.
    This does not produce realistic particle masses, and we
    do not claim it does.
    The conductor--mass relationship, if it exists, must be
    more subtle than a direct identification.

  \item \textbf{$\GL(n)$ vs.\ $\SU(n)$.}
    The passage from $\GL(n)$ (automorphic side) to $\SU(n)$
    (gauge side) involves taking compact forms and modding out
    centers.
    This step is physically motivated but not derived from
    first principles within the arithmetic framework.

  \item \textbf{The $1.3\%$ discrepancy.}
    The predicted $\OmL = 2\pi/9 \approx 0.698$ is $\sim 1.3\%$
    above the Planck central value $0.689$.
    While within $2\sigma$, this may indicate missing corrections
    (e.g., from higher-order arithmetic terms).

  \item \textbf{No dynamics.}
    The framework determines the gauge group and spacetime
    dimension but does not derive the dynamics (Lagrangian)
    of the Standard Model.
    It is a ``kinematic'' framework, not a dynamical one.
\end{enumerate}

%======================================================================
\section{Conclusion}\label{sec:conclusion}
%======================================================================

The \'etale cohomological dimension $\cd(\Spec(\Z))=3$ is the
``master number'' of physics.
From this single arithmetic invariant, we derive:
\begin{itemize}
  \item Spacetime dimension: $d = \cd + 1 = 4$.
  \item Standard Model gauge group:
    $\U(1)\times\SU(2)\times\SU(3)$ from $\GL(1)\times\GL(2)\times\GL(3)$
    with $n\le\cd=3$.
  \item Gauge boson count: $1+3+8 = 12 = 2/|B_2|$.
  \item Dark energy:
    $\OmL = (4/3)\times 2\pi\times(1/12) = 2\pi/9$.
\end{itemize}

The dark energy decomposes into three factors---gravitational,
archimedean, and arithmetic---reflecting the Arakelov structure
of $\Spec(\Z)$.
The archimedean factor ($2\pi$) explains why dark energy exists
in our universe: because our universe is built on real numbers,
not $p$-adic numbers.

The Eisenstein/cusp decomposition of modular forms provides
a bridge between the vacuum/matter structure of physics and
the arithmetic of $\Spec(\Z)$, with Eisenstein series encoding
the non-perturbative vacuum (left wing) and cusp forms encoding
perturbative matter (right wing).

We believe that the path from $\cd(\Spec(\Z))=3$ to the full
structure of the Standard Model is one of the most promising
directions in mathematical physics, bridging arithmetic geometry
and fundamental physics in a way that neither field could
achieve alone.

%======================================================================
\begin{thebibliography}{99}

\bibitem{ArtinVerdier}
M.~Artin and J.-L.~Verdier,
``Seminar on \'Etale Cohomology of Number Fields,''
Woods Hole, 1964.

\bibitem{Milne_ADT}
J.~S.~Milne,
\textit{Arithmetic Duality Theorems}, 2nd ed.,
BookSurge, 2006.

\bibitem{Milne_EC}
J.~S.~Milne,
\textit{\'Etale Cohomology},
Princeton University Press, 1980.

\bibitem{Morishita}
M.~Morishita,
\textit{Knots and Primes: An Introduction to Arithmetic Topology},
Springer, 2012.

\bibitem{Langlands}
R.~P.~Langlands,
``Problems in the theory of automorphic forms,''
\textit{Lectures in Modern Analysis and Applications III},
Lecture Notes in Math.\ \textbf{170}, Springer, 1970, pp.~18--61.

\bibitem{Wiles}
A.~Wiles,
``Modular elliptic curves and Fermat's Last Theorem,''
\textit{Ann.\ of Math.}\ \textbf{141} (1995), 443--551.

\bibitem{BCDT}
C.~Breuil, B.~Conrad, F.~Diamond, and R.~Taylor,
``On the modularity of elliptic curves over $\Q$,''
\textit{J.\ Amer.\ Math.\ Soc.}\ \textbf{14} (2001), 843--939.

\bibitem{Scholze}
P.~Scholze,
``On torsion in the cohomology of locally symmetric varieties,''
\textit{Ann.\ of Math.}\ \textbf{182} (2015), 945--1066.

\bibitem{Arakelov}
S.~J.~Arakelov,
``Intersection theory of divisors on an arithmetic surface,''
\textit{Izv.\ Akad.\ Nauk SSSR Ser.\ Mat.}\ \textbf{38} (1974), 1179--1192.

\bibitem{Neukirch}
J.~Neukirch, A.~Schmidt, and K.~Wingberg,
\textit{Cohomology of Number Fields}, 2nd ed.,
Springer, 2008.

\bibitem{WB_spacetime}
Wright Brothers Collaboration,
``Spacetime dimension from $\cd(\Spec(\Z))=3$,''
2025.

\bibitem{WB_darkenergy}
Wright Brothers Collaboration,
``Dark energy from the Riemann zeta function,''
2025.

\bibitem{Planck2018}
Planck Collaboration,
``Planck 2018 results. VI. Cosmological parameters,''
\textit{Astron.\ Astrophys.}\ \textbf{641} (2020), A6.

\bibitem{HarrisTaylor}
M.~Harris and R.~Taylor,
\textit{The Geometry and Cohomology of Some Simple Shimura Varieties},
Ann.\ of Math.\ Studies \textbf{151}, Princeton Univ.\ Press, 2001.

\bibitem{Clozel}
L.~Clozel,
``Motifs et formes automorphes,''
\textit{Automorphic Forms, Shimura Varieties, and $L$-functions},
Academic Press, 1990.

\end{thebibliography}

\end{document}
